\chapter{Dataset and Data Preparation}
\label{ch:dataset}
% Chapter 3: Dataset and Data Preparation (10-12 pages)

%
% Section: Data Sources
%
\section{Data Sources}
\label{sec:dataset:sources}

\subsection{MultiCaRe Dataset Overview}
\label{subsec:dataset:sources:multicare}

\subsection{PubMed Central Case Reports}
\label{subsec:dataset:sources:pmc}

\subsection{Data Extraction Pipeline}
\label{subsec:dataset:sources:extraction}

%
% Section: Leishmaniasis Case Curation
%
\section{Leishmaniasis Case Curation}
\label{sec:dataset:curation}

\subsection{Case Selection Criteria}
\label{subsec:dataset:curation:criteria}

\subsection{Verification Pipeline}
\label{subsec:dataset:curation:verification}
% 406 → 143 verified cases

\subsection{Diagnosis Type Classification}
\label{subsec:dataset:curation:diagnosis}
% CL, VL, MCL, PKDL

\subsection{Species Identification}
\label{subsec:dataset:curation:species}

%
% Section: Image Processing
%
\section{Image Processing}
\label{sec:dataset:images}

\subsection{Clinical Image Extraction}
\label{subsec:dataset:images:extraction}

\subsection{Caption Alignment}
\label{subsec:dataset:images:captions}

\subsection{Image Quality Filtering}
\label{subsec:dataset:images:quality}

%
% Section: Train/Test Split Design
%
\section{Train/Test Split Design}
\label{sec:dataset:split}

\subsection{Article-based Splitting}
\label{subsec:dataset:split:article}
% No data leakage

\subsection{Split Statistics}
\label{subsec:dataset:split:statistics}
% 114 train / 29 test

\subsection{Disjointness Verification}
\label{subsec:dataset:split:verification}

%
% Section: Query Generation
%
\section{Query Generation}
\label{sec:dataset:queries}

\subsection{Query Types}
\label{subsec:dataset:queries:types}
% Q1 (text-only), Q1+Q3 (multimodal), Q3 (image-only)

\subsection{Symptom Extraction with Leak Pattern Filtering}
\label{subsec:dataset:queries:extraction}

\subsection{Query Template Design}
\label{subsec:dataset:queries:templates}

%
% Section: Ground Truth Generation
%
\section{Ground Truth (QRels) Generation}
\label{sec:dataset:qrels}

\subsection{Relevance Grading}
\label{subsec:dataset:qrels:grading}
% 0-3 scale

\subsection{Diagnosis Type Matching}
\label{subsec:dataset:qrels:diagnosis}

\subsection{Species-level Matching}
\label{subsec:dataset:qrels:species}
